\section{Mechanism Classes and Their Physical Signatures}\label{sec:signatures}
% Budget 2.5 pages. Table 2.
% Table 2: mechanism classes and their physical signatures (Plan §9; plan/CHANGELOG.md C6).
% Cell sources: paper/tables/table2_NOTES.md. Needs \usepackage{booktabs,tabularx} in main.tex.
\begin{table}[t]
\caption{Predicted response of each mechanism class to each probe of a frozen deadlock. ``---'': the probe does not test that class.}
\label{tab:signatures}
\centering
\scriptsize
\setlength{\tabcolsep}{2pt}
\renewcommand{\arraystretch}{1.15}
\begin{tabularx}{\textwidth}{@{}l*{7}{>{\raggedright\arraybackslash}X}@{}}
\toprule
Class
 & Quench
 & Replay, $10^{-9}$\,m
 & $\Delta t$, per-step noise
 & $\Delta t$, $\sqrt{\Delta t}$ noise
 & Rest-drive push, $+$ vs $-$
 & Tilt $b\to-b$
 & Corrected numerics \\
\midrule
P
 & held
 & held
 & held
 & held
 & unchanged
 & unchanged
 & persists \\
L
 & frees, deterministic
 & identical replays
 & no $\Delta t$ dependence
 & no $\Delta t$ dependence
 & $+$ frees, $-$ does not
 & ---
 & persists \\
T
 & frees at random, exponential times
 & replays disagree
 & rate falls as $\Delta t$ grows
 & no $\Delta t$ dependence
 & ---
 & ---
 & persists \\
D
 & frees only beyond $b_f$
 & identical replays
 & no $\Delta t$ dependence
 & no $\Delta t$ dependence
 & ---
 & first out switches side
 & persists \\
N
 & frees or holds, depending on solver details
 & outcome changes under $10^{-9}$\,m
 & ---
 & ---
 & ---
 & ---
 & effect disappears \\
\bottomrule
\end{tabularx}
\end{table}

%
A deadlock snapshot is a rest state of Eq.~\eqref{eq:force}: every agent is still and the forces on it balance. Removing the blocker frees a lone keystone; at a second trap another agent takes its place and is held; in a group standoff several agents stay held. An intervention can act on a deadlock in five physically distinct ways. \textbf{P, prevent:} it acts only while the crowd moves, changing how often deadlocks form but leaving a frozen one as it is. \textbf{L, reshape:} it adds a deterministic force that persists at rest, so the fixed point moves or disappears. \textbf{T, heat:} it adds fluctuations; the fixed point survives and is left over a barrier $\Delta U$ at a rate set by $\Delta U/k_BT_{\rm eff}$. \textbf{D, steer:} a finite directional bias breaks a symmetry; a strict minimum survives small biases and disappears at a bias $b_f$. \textbf{N, glitch:} the outcome changes under a perturbation far below the model's physical scales, or the effect vanishes under corrected numerics. Each class predicts a response (Table~\ref{tab:signatures}) to probes of the same frozen state: a quench under the intervention, a replay perturbed by $10^{-9}$\,m, a time-step change with per-step or $\sqrt{\Delta t}$-scaled noise, a constant push equal to ORCA's rest drive (the rest-drive push, $+$ towards the door, $-$ away), reversal of a lateral tilt, and corrected numerics. The probes are paired by snapshot, and none is used to tune the control. The Jaccard index $J$ of the outcome sets, the time to free, the final-state class and stall survival are descriptive signatures beyond Table~\ref{tab:signatures}; $J$ separates resolvers only at the pair standoff, where both are deterministic.

\paragraph{Attribution protocol.}
(1)~Choose a resolver credited with resolving deadlocks and a control from a different class. (2)~Tune the control's single parameter to the resolver's rate on the same deadlocks (freed fraction for snapshots, completion fraction for whole runs); this is calibration, not evidence. (3)~Apply the probes of Table~\ref{tab:signatures} to both on a common library of snapshots. A probe separates two interventions when their 95\% intervals for that signature exclude each other, or when one shows the predicted effect and the other its absence; replay and handedness are exact counts. (4)~Assign each intervention to the class whose responses it shows. The resolver and the control share a mechanism only if no probe separates them. The test is informative because the resolver could have shown any class's signature, and it can confirm: at the lone keystone the rest-drive push frees every blocker ORCA frees, replays exactly and, reversed, frees none (Sect.~\ref{sec:att-lone}).

\subsection{Terms That Vanish at Rest}
Stack the $n$ remaining agents into $\mathbf x,\mathbf v\in\mathbb R^{2n}$, freeze the densities at their snapshot values and split the force as
\begin{equation}\label{eq:split}
\mathbf F(\mathbf x,\mathbf v)=\mathbf F_0(\mathbf x)-\mathsf D\,\mathbf v+\mathbf G(\mathbf x,\mathbf v),
\end{equation}
with the rest-state field $\mathbf F_0(\mathbf x)=\mathbf F(\mathbf x,\mathbf 0)$, the damping $\mathsf D={\rm diag}(m_i/\tau_i)$ of Eq.~\eqref{eq:des}, and the rest $\mathbf G$, with $\mathbf G(\mathbf x,\mathbf 0)=\mathbf 0$.

\begin{proposition}[rest-state lemma]\label{prop:rest}
Suppose $\mathbf v\cdot\mathbf G(\mathbf x,\mathbf v)\le0$ for all $|\mathbf v|<\epsilon$. Then (i) the fixed points of $\mathsf M\dot{\mathbf v}=\mathbf F$, $\dot{\mathbf x}=\mathbf v$, with $\mathsf M$ the mass matrix, are exactly the states $(\mathbf x^\ast,\mathbf 0)$ with $\mathbf F_0(\mathbf x^\ast)=\mathbf 0$; and (ii) if $\mathbf F_0=-\nabla U$, every strict local minimum of $U$ is a Lyapunov-stable fixed point.
\end{proposition}
\begin{proof}
(i) A fixed point has $\mathbf v=\mathbf 0$ and $\mathbf F(\mathbf x^\ast,\mathbf 0)=\mathbf F_0(\mathbf x^\ast)=\mathbf 0$; conversely such a state does not move. (ii) Let $E=\tfrac12\mathbf v^{\top}\mathsf M\mathbf v+U(\mathbf x)$. Along trajectories $\dot E=\mathbf v\cdot(\mathbf F+\nabla U)=-\mathbf v^{\top}\mathsf D\,\mathbf v+\mathbf v\cdot\mathbf G\le0$ while $|\mathbf v|<\epsilon$. Near a strict minimum, small sublevel sets of $E$ around $(\mathbf x^\ast,\mathbf 0)$ trap trajectories.
\end{proof}

The hypotheses hold term by term. The frictions of Eqs.~\eqref{eq:sfm}--\eqref{eq:wall} vanish at rest and dissipate: their power is $-\kappa\,g\,[(\mathbf v_i-\mathbf v_j)\cdot\mathbf t_{ij}]^2\le0$ per pair and $-\kappa\,g\,(\mathbf v_i\cdot\mathbf t_{iW})^2\le0$ per wall. The TTC force (Eq.~\eqref{eq:ttc}) is zero below relative speeds of $10^{-4}$\,m/s. With the densities frozen, the C1 rest field is conservative,
\begin{equation}\label{eq:U}
U(\mathbf x)=\sum_i c_i\,|\mathbf x_i-\mathbf g|+\sum_{i<j}\phi_{ij}(d_{ij})+\sum_{i,W}\phi_{iW}(d_{iW}),\qquad c_i=\frac{m_iv_{{\rm eff},i}}{\tau_i},
\end{equation}
with $\phi$ the potentials of the exponential and compression terms; in the code $\max|\mathbf F_0+\nabla U|\le\ResidCollMax$\,N over the \ResidBarN\ barrier and $\ResidPairMax$\,N over the \ResidPairN\ pair snapshots. Hence C2 has exactly C1's fixed points and stability, and neither TTC nor friction can resolve a deadlock (class P at most). ORCA's rest drive does not vanish at rest; it enters $\mathbf F_0$ and reshapes the landscape (a candidate for class L). The lemma covers the noise-free model only, and part (ii) covers C1 and C2 but not ORCA's velocity-dependent constraint part.

\subsection{Effective Temperature of Per-Step Noise}\label{sec:teff}
Near rest and away from contact, one velocity component of agent $i$ under the corrected update and the damping of Eq.~\eqref{eq:des} is the autoregressive process $v_{n+1}=a\,v_n+\xi_n$, $a=1-\Delta t/\tau_i$, $\langle\xi_n\xi_{n'}\rangle=\sigma^2\delta_{nn'}$, with stationary variance $\sigma^2/(1-a^2)$. Setting $m_i\langle v^2\rangle=k_BT_{\rm eff}$ gives
\begin{equation}\label{eq:teff}
k_BT_{\rm eff}=\frac{m_i\,\sigma^2}{1-(1-\Delta t/\tau_i)^2}=\frac{m_i\,\sigma^2\tau_i}{2\,\Delta t}\left(1-\frac{\Delta t}{2\tau_i}\right)^{-1}.
\end{equation}
A restoring curvature $K$ adds terms of order $(\omega\Delta t)^2$, $\omega^2=K/m_i$, negligible in our snapshots. At fixed per-step $\sigma$, $k_BT_{\rm eff}\propto1/\Delta t$ up to the factor $(1-\Delta t/2\tau_i)^{-1}$, which $\sqrt{\Delta t}$-scaled noise alone leaves. Eq.~\eqref{eq:teff} assumes frozen densities, linear friction and no contact or speed clipping; through $\tau_i$ the temperature varies up to 1.5-fold between agents (Table~\ref{tab:params}), and $\sigma$ is set from the blocker's $\tau_i$. The blocker's realised velocity variance over Eq.~\eqref{eq:teff} is \TthreeRatioMean\ in dedicated runs (median \TthreeCellsMedian\ over the \TthreeCells\ noise cells of the escape experiments).

If escape is activated, $\lambda\propto\exp(-\Delta U/k_BT_{\rm eff})$ and the log rate ratio between two time steps has no free parameter,
\begin{equation}\label{eq:ell}
\ell\equiv\ln\frac{\lambda(\Delta t_1)}{\lambda(\Delta t_2)}=\Delta U\left(\frac{1}{k_BT_{\rm eff}(\Delta t_2)}-\frac{1}{k_BT_{\rm eff}(\Delta t_1)}\right),
\end{equation}
positive for $\Delta t_1<\Delta t_2$ at fixed per-step $\sigma$, near zero for $\sqrt{\Delta t}$-scaled noise, and absent for a deterministic resolver.

\subsection{Barriers from the String Method}\label{sec:barrier}
The barrier in Eq.~\eqref{eq:ell} belongs to the many-agent landscape of Eq.~\eqref{eq:U}. We compute it with the string method~\cite{e2007} on all mobile coordinates, from the snapshot minimum to an exit state (blocker past the door plane), with all residual agents mobile (full), the neighbours pinned (pinned), or the blocker alone (alone). A barrier is used only if the string has converged (uphill perpendicular force below 5\,N, climbing image stationary at the path maximum, under 1\% change on doubling the images), ends in an exit, and, when alone, matches a direct two-dimensional path. Kramers theory predicts~\cite{kramers1940,hanggi1990}: escape times are exponential (Weibull shape $k=1$)~\cite{weibull1951}; $\ln\lambda$ is linear in $1/k_BT_{\rm eff}$ with slope $-\Delta U$ for the condition's own barrier; and, since the neighbours can move out of the way, the full barrier lies below the alone barrier.

